Seluruh pengujian pada penelitian ini dilaksanakan pada lingkungan simulasi Gazebo~11 yang terintegrasi dengan ROS Noetic sesuai arsitektur sistem yang telah dijelaskan pada Bab~III. Proses eksperimen dilakukan untuk mengevaluasi pengaruh integrasi kamera termal, fusi trav\-ersability LiDAR--termal, serta pengendali fuzzy Mamdani terhadap performa navigasi robot bergerak pada berbagai kondisi lingkungan. Fokus evaluasi diarahkan pada kemampuan sistem dalam mempertahankan margin keselamatan navigasi, mendeteksi obstacle pada area \textit{blind-spot} LiDAR, serta menjaga efisiensi pergerakan robot menuju tujuan.

Pelaksanaan eksperimen mengikuti prosedur pengujian yang telah dijelaskan pada Subbab~\ref{sec:eksperimen}. Setiap skenario dijalankan sebanyak sepuluh kali percobaan independen ($n = 10$) menggunakan konfigurasi awal yang identik. Dengan delapan skenario pengujian yang terdiri atas empat skenario \textit{baseline} dan empat skenario \textit{fuzzy}, total diperoleh 80 \textit{trial} yang digunakan sebagai dasar analisis pada bab ini. Seluruh data eksperimen direkam secara otomatis menggunakan node \texttt{experiment\_logger.py} dan disimpan ke dalam direktori \texttt{experiment\_data} untuk selanjutnya diproses pada tahap analisis statistik.

Kedelapan skenario pengujian disusun dalam empat pasangan perbandingan yang mewakili kondisi lingkungan yang sama namun menggunakan konfigurasi kendali yang berbeda. Pas\-angan S1 dan S5 merepresentasikan lingkungan terbuka tanpa obstacle. Pasangan S2 dan S6 merepresentasikan lingkungan dengan satu obstacle manusia statis. Pasangan S3 dan S7 merepresentasikan lingkungan dengan beberapa obstacle statis yang mencakup keberadaan obstacle berprofil rendah berupa kucing pada area \textit{blind-spot} LiDAR. Sementara itu, pasangan S4 dan S8 merepresentasikan lingkungan dengan obstacle dinamis yang bergerak melintasi jalur robot serta keberadaan obstacle kucing pada area \textit{blind-spot}. Konfigurasi tersebut dipilih untuk mengevaluasi performa sistem pada tingkat kompleksitas lingkungan yang semakin meningkat.

Pada seluruh skenario, konfigurasi sensor yang digunakan tetap sama, yaitu sensor LiDAR 2D dan kamera termal yang aktif secara bersamaan. Kedua sensor diproses melalui \textit{pipeline} yang identik sehingga menghasilkan informasi traversability berbasis LiDAR, traversability berbasis termal, serta traversability efektif hasil fusi sensor. Informasi tersebut kemudian digunakan oleh APFPlanner untuk menghasilkan perintah navigasi robot. Perbedaan utama antara kelompok \textit{baseline} (S1--S4) dan kelompok \textit{fuzzy} (S5--S8) terletak pada penggunaan pengendali fuzzy Mamdani sebagai mekanisme adaptasi parameter navigasi. Pada konfigurasi \textit{baseline}, parameter APF menggunakan nilai tetap, sedangkan pada konfigurasi \textit{fuzzy}, parameter \textit{speed scale} dan \textit{repulsive gain} disesuaikan secara adaptif berdasarkan nilai jarak obstacle terdekat dan traversability efektif sebagaimana dijelaskan pada Subbab~\ref{sec:fuzzy_Aapf}.

Data primer yang diperoleh dari setiap \textit{trial} meliputi status keberhasilan navigasi, jumlah tabrakan, nilai minimum clearance, waktu tempuh, panjang lintasan, kecepatan linear rata-rata, dan kecepatan angular rata-rata. Seluruh data tersebut direkam secara \textit{real-time} selama eksperimen berlangsung dan kemudian diproses menggunakan skrip analisis untuk menghasilkan statistik deskriptif serta hasil uji inferensial. Analisis statistik dilakukan menggunakan uji Mann--Whitney U dan \textit{rank-biserial correlation} sesuai metode yang telah dijelaskan pada Subbab~\ref{sec:evaluasi}.

Penyajian hasil pada bab ini mengikuti urutan tujuan evaluasi yang telah ditetapkan pada Subbab~\ref{sec:evaluasi}. Pembahasan diawali dengan penyajian hasil deskriptif pada masing-masing skenario untuk memberikan gambaran umum performa sistem. Selanjutnya dilakukan evaluasi tingkat keberhasilan navigasi dan ketiadaan tabrakan pada seluruh skenario pengujian. Analisis kemudian difokuskan pada \textit{minimum clearance} sebagai metrik utama yang digunakan untuk menilai peningkatan keselamatan navigasi akibat integrasi kamera termal dan pengendali fuzzy. Setelah itu dibahas \textit{trade-off} yang muncul terhadap efisiensi navigasi melalui analisis waktu tempuh dan panjang lintasan, serta kehalusan pergerakan robot berdasarkan karakteristik kecepatan angular. Untuk memperkuat temuan kuantitatif, beberapa visualisasi trajektori robot juga disajikan pada skenario yang memiliki tingkat kompleksitas tertinggi. Seluruh hasil pengujian kemudian dirangkum dan dibahas secara terintegrasi untuk menjawab tujuan penelitian yang telah dirumuskan pada Bab~I.
